\subsection{Compatibilité entre les représentations dodécaphasique et analytique d'alpha}
\label{subsec:alpha-limite-dos-vias}

La compatibilité s'établit à toute profondeur. La représentation
dodécaphasique réversible et la représentation analytique de la
coordonnée d'alpha
\(E_{21/22}\) sont deux représentations internes de codomaines distincts,
produites à partir du même état enrichi APP--TRIT--TPK. Le jet
\(E_{21/22}^{(9)}\) est un témoin fini exact, mais ne détermine pas à lui seul
la queue. La complétion construite ensuite ne se déduit pas du jet
seul : elle reçoit la coordonnée partagée de l'état et exprime son image
analytique par une récurrence explicite. Elle ne constitue donc pas une
sélection causale indépendante ; elle constitue une représentation analytique
typée du même point engendré.

Cette séparation conserve l'ordre causal
\[
 \mathrm{APP}\longrightarrow\mathrm{TRIT}\longrightarrow\mathrm{TPK}
 \longrightarrow\mathcal X_{\alpha}^{\mathrm{enr}}
 \longrightarrow (P,E,\Phi,U_{12}^{\mathrm{sgn}},K)
 \longrightarrow\alpha_{\mathrm{HMT}}.
\]
Les coordonnées \(P,E,\Phi\) et le registre \(K\) sont des sorties antérieures du
même état enrichi. Leur utilisation dans la clôture d'alpha est donc une
composition interne ultérieure et non une injection de constantes
conventionnelles.

Pour \(r\in6\mathbb N\), \(r\ge36\), soit
\(\widehat{\mathcal X}^{\mathrm{enr},\alpha}_r\) le secteur des états
tronqués qui prolongent le germe interne d'alpha fixé dans \(R_{36}\). On
écrit
\[
 x_r^\alpha=(\widetilde x_r,\chi_\alpha,
              \varepsilon_r^\alpha,N_r^\alpha),
\]
où \(\chi_\alpha\) est le caractère interne du secteur, et non la valeur scalaire
exprimée. Chaque état conserve reste, quotient, retenue, feuillet, orientation,
phase, route, cylindre, frontière, signature, registre, mémoire, survie,
provenance et profondeur. Si
\[
 d_r^\alpha=\operatorname{Dig}_{729}(\varepsilon_r^\alpha),\qquad
 \varepsilon_{r+6}^\alpha=729\varepsilon_r^\alpha-d_r^\alpha,
 \qquad N_{r+6}^\alpha=729N_r^\alpha+d_r^\alpha,
\]
la transition effective est
\[
 \mathcal U_r^\alpha=
 \operatorname{Upd}_r^\alpha\circ
 \operatorname{Tra}_r^\alpha\circ
 \operatorname{Sel}_r^\alpha,
 \qquad
 \operatorname{Ext}^{\alpha,\to}_{r+6,r}
 =\operatorname{graph}(\mathcal U_r^\alpha).
\]
La fonctionnalité de ces extensions, une fois le germe \(R_{36}\) fixé,
produit une unique section compatible : la phase revient tandis que la mémoire et
la profondeur avancent. Le domaine global des deux déterminations est
\begin{equation}
 \mathcal X_\alpha^{\mathrm{enr}}:=\widehat\Omega_\alpha
 :=\varprojlim_{r\ge36}
 \bigl(\widehat{\mathcal X}^{\mathrm{enr},\alpha}_r,
       \operatorname{Ext}^{\alpha,\to}_{r+6,r}\bigr).
 \label{eq:alpha-dominio-enriquecido-completo}
\end{equation}

\subsubsection{La clôture dodécaphasique comme section complète}

Pour \(N\geq1\), posons
\[
 \mathcal W_N:=\{0,\ldots,999\}^{N},
 \qquad
 \mathsf T_{\alpha,N}(\mathsf s):=(A_1,\ldots,A_N),
 \qquad \mathsf s\in\mathcal X_{\alpha}^{\mathrm{enr}},
\]
où les blocs \(A_j\) s'obtiennent au moyen de
\eqref{eq:alpha-recurrencia} avec la frontière compatible
\(c_{N+1}^{\infty}=\lfloor R_{N+1}(Z)\rfloor\). Si \(N'\geq N\), soit
\(\rho_{N',N}:\mathcal W_{N'}\to\mathcal W_N\) la restriction aux
\(N\) premiers blocs.

\begin{theorem}[Sortie dodécaphasique complète et compatibilité projective]
\label{thm:alpha-salida-dodecafasica-completa}
La collection \((\mathsf T_{\alpha,N})_{N\geq1}\) est compatible :
\begin{equation}
 \rho_{N',N}\circ\mathsf T_{\alpha,N'}
 =\mathsf T_{\alpha,N}
 \qquad(N'\geq N).
 \label{eq:alpha-via-a-proyectiva-exacta}
\end{equation}
Ses cylindres canoniques
\begin{equation}
 C_N^A(\mathsf s):=
 \left[
  \sum_{j=1}^{N}A_j1000^{-j},
  \sum_{j=1}^{N}A_j1000^{-j}+1000^{-N}
 \right)\cap I_\alpha
 \label{eq:alpha-cilindros-via-a-exactos}
\end{equation}
sont emboîtés, ont un diamètre au plus égal à \(1000^{-N}\) et satisfont
\begin{equation}
 \bigcap_{N\geq1}C_N^A(\mathsf s)
 =\{\alpha_A(\mathsf s)\},
 \qquad
 \alpha_A
 =\pi_{\mathrm{HMT}}+e_{\mathrm{HMT}}-\varphi_{\mathrm{HMT}}
  -4-\kappa_{\mathrm{per}}.
 \label{eq:alpha-salida-a-completa}
\end{equation}
Par conséquent, la première voie exprime une unique section infinie, que
nous notons \(\alpha_{\mathrm{HMT}}:=\alpha_A\).
\end{theorem}

\begin{proof}
La proposition~\ref{prop:alpha-normalizacion} fournit l'unique développement
canonique et l'unique suite bornée de retenues. L'identité locale
\eqref{eq:alpha-identidad-local} et sa sommation télescopique
\eqref{eq:alpha-telescopia} conservent la frontière distante ; la proposition
\ref{ampl-alfa-prop-frontera} prouve que restreindre une exécution prolongée
restitue exactement le préfixe antérieur. Cela donne
\eqref{eq:alpha-via-a-proyectiva-exacta}. L'appartenance de la somme complète
à chaque cylindre découle de l'identité des queues, et
\(\operatorname{diam}C_N^A\leq1000^{-N}\to0\) ; l'intersection contient un
seul point. Enfin, \eqref{eq:alpha-identidad-completa} identifie ce
point à l'expression de \eqref{eq:alpha-salida-a-completa}. Tous ses
coefficients et signes sont fixés avant la représentation d'alpha.
\end{proof}

La fenêtre \(\alpha_{12}^{(0)}\), fermée avec \(c_{13}=0\), reste un
calcul fini exact ; elle ne doit pas être confondue avec la restriction de la section
complète, dont la frontière compatible satisfait \(c_{13}^{\infty}=-1\). En
particulier,
\begin{equation}
 \rho_{36,12}(\alpha_{A,36})=\alpha_{A,12}
 \label{eq:alpha-rho-36-12-proyectiva}
\end{equation}
est une égalité de mots compatibles, et non une égalité entre des mots de
longueurs différentes ni une identification de la fenêtre isolée à la
limite.

\subsubsection{Caractérisation analytique exacte d'ordre neuf}

Soit
\[
 \mathbb K:=\mathbb Q\!\left(\pi_{\mathrm{HMT}},
                  \log_{10}\frac{10}{9}\right),
 \qquad
 E_9(x):=E_{21/22}^{(9)}(x)=P_9(x)-\delta .
\]
Les équations \eqref{eq:alpha-firma-coeficientes} et
\eqref{eq:alpha-t79} construisent \(P_9=P_6+T_{7:9}\) à partir de la signature de l'
état et de la résolvante nilpotente ; la
proposition~\ref{prop:alpha-raiz-jet} démontre que \(E_9\) possède une unique
racine simple \(\xi_9\in I_\alpha\), isolée par l'intervalle rationnel
\eqref{eq:alpha-intervalo-jet}.

La bidirectionnalité finie peut s'écrire sans introduire de queue
analytique. Définissons
\begin{align}
 B_9(x):={}&\delta+\frac74x^2-\frac{x^4}{20}
 +\frac{2x^5}{21}+\frac{x^6}{46}+\frac{x^7}{120}
 -\frac{x^8}{45}-\frac{2x^9}{495},
 \label{eq:alpha-b9-finito}\\
 q_x(T):={}&xT^2-B_9(x)T+x^3,\qquad x\in I_\alpha=(0,10^{-2}),
 \label{eq:alpha-q-involucion-finita}\\
 J_x:{\mathbb R}^{\times}\longrightarrow{\mathbb R}^{\times},
 \qquad J_x(T):={}&\frac{x^2}{T}.
\end{align}

\begin{proposition}[Équivalence quadratique et involution finies]
\label{prop:alpha-involucion-finita-exacta}
Pour tout \(x\in I_\alpha\) et tout \(T\in\mathbb R^\times\),
\begin{align}
 \frac{q_x(2\pi_{\mathrm{HMT}})}{2\pi_{\mathrm{HMT}}}
 &=P_9(x)-\delta=E_9(x),
 \label{eq:alpha-q-equivale-p9}\\
 J_x(J_x(T))&=T,
 \qquad
 T^2q_x(J_x(T))=x^2q_x(T).
 \label{eq:alpha-j-involucion-exacta}
\end{align}
Soit \(Z_x=\{T\in\mathbb R^\times:q_x(T)=0\}\). Par conséquent,
\(J_x|_{Z_x}:Z_x\to Z_x\) est involutive. Lorsque \(E_9(x)=0\),
elle échange les racines \(2\pi_{\mathrm{HMT}}\) et
\(x^2/(2\pi_{\mathrm{HMT}})\). Il s'agit d'une involution exacte sur la fibre
scalaire du jet d'ordre neuf ; elle n'agit pas sur l'état enrichi.
\end{proposition}

\begin{proof}
Diviser \(q_x(T)\) par \(T\) et substituer \(T=2\pi_{\mathrm{HMT}}\) reproduit terme à terme la définition de \(P_9(x)-\delta\). La première identité de \eqref{eq:alpha-j-involucion-exacta} est immédiate. Pour la seconde,
\[
 T^2q_x(x^2/T)
 =x^5-B_9(x)x^2T+x^3T^2=x^2q_x(T).
\]
Ainsi, \(J_x\) préserve l'ensemble des racines et, par la formule de Viète,
échange les deux racines dont le produit est \(x^2\).
\end{proof}

Le témoin décimal fini conserve alors sa force exacte :
\begin{equation}
 \operatorname{Tr}_9(\xi_9^{-1})
 =\operatorname{Tr}_9\!\left((\alpha_{12}^{(0)})^{-1}\right)
 =137.035999084.
 \label{eq:alpha-testigo-finito-nueve}
\end{equation}
L'égalité \eqref{eq:alpha-testigo-finito-nueve} certifie la caractérisation finie et l'involution précédente. Elle ne transforme pas la racine d'un polynôme tronqué en la racine de tout prolongement compatible.

\subsubsection{Le jet ne détermine pas à lui seul une prolongation analytique}

Soit
\[
 j_9:\mathbb K[[x]]\longrightarrow \mathbb K[x]/(x^{10})
\]
la troncature, et pour \(h\in\mathbb K[x]\), posons
\begin{equation}
 E_h(x):=E_9(x)+x^{10}h(x).
 \label{eq:alpha-familia-prolongaciones-mismo-jet}
\end{equation}
Tous ces polynômes satisfont \(j_9(E_h)=j_9(E_9)\).

\begin{theorem}[Non-unicité de la racine à partir du jet]
\label{thm:alpha-no-unicidad-prolongacion-jet}
Il n'existe pas d'application
\[
 \mathfrak r_9:\mathbb K[x]/(x^{10})\longrightarrow I_\alpha
\]
qui attribue, à partir du jet seul, la racine positive de toute
prolongation polynomiale de \(E_9\) possédant une unique racine dans
\(I_\alpha\). En particulier, \(E_0\) et \(E_{1/729}\) possèdent le même jet et
des racines positives distinctes.
\end{theorem}

\begin{proof}
La borne démontrée dans la preuve de
\ref{prop:alpha-raiz-jet} donne \(E_9'(x)>6.24\) sur
\([0,10^{-2}]\). Pour \(t\geq0\), écrivons
\(E_t(x)=E_9(x)+tx^{10}\). Alors
\[
 E_t'(x)=E_9'(x)+10tx^9>6.24,
\]
tandis que \(E_t(0)=-\delta<0\) et \(E_t(10^{-2})>E_9(10^{-2})>0\). Chaque \(E_t\) possède donc une unique racine positive \(\eta(t)\in I_\alpha\). Pour \(t=1/729\),
\[
 E_{1/729}(\xi_9)=\frac{\xi_9^{10}}{729}>0;
\]
comme \(E_{1/729}\) est strictement croissant et ne change de signe qu'une seule
fois, sa racine satisfait \(\eta(1/729)<\xi_9=\eta(0)\). Cependant,
\(j_9(E_{1/729})=j_9(E_0)\). Une application qui factoriserait la racine par
\(j_9\) devrait attribuer le même point aux deux polynômes, ce qui est contradictoire.
\end{proof}

La variation n'est pas seulement existentielle. Le théorème des fonctions implicites s'applique puisque \(\partial_xE_t(\eta(t))>6.24\), et fournit
\begin{equation}
 \eta'(t)=-
 \frac{\eta(t)^{10}}
 {E_9'(\eta(t))+10t\eta(t)^9}<0.
 \label{eq:alpha-variacion-raiz-prolongacion}
\end{equation}
Il est ainsi prouvé que les coefficients postérieurs au degré neuf modifient
la racine même lorsqu’ils laissent invariant tout le jet construit jusqu’à ce degré.
Cette conclusion ne nie pas la queue produite par l’état enrichi : elle prouve
uniquement que la racine ne se factorise pas par la troncature nue
\(j_9\). Choisir une queue \(h\) à partir d’une racine conventionnelle externe
serait circulaire. La complétion qui suit ne consulte ni cette racine ni une table
métrologique : elle évalue directement les coordonnées régionales et terminales de
l’état. Comme cette évaluation coïncide algébriquement avec la représentation
entière, le résultat est une représentation corrélée, non une affirmation
d’indépendance entre deux sélecteurs.

% Representación analítica correlacionada de alfa: acción, convergencia e
% intervalos.
\subsubsection{Représentation analytique de la coordonnée d'alpha}
\label{subsec:alpha-lector-completo-rev08}

Le contrôle du jet précédent oblige à exhiber la queue : il ne suffit pas de nommer un
lecteur complet. Soit \(\mathsf s\in\mathcal X_\alpha^{\mathrm{enr}}\). Avant
la normalisation décimale, le même état exprime déjà la précoordonnée
\begin{equation}
 a(\mathsf s):=p(\mathsf s)+e_0(\mathsf s)-f(\mathsf s)
                 -\kappa_{\mathrm{per}}(\mathsf s)\in I_\alpha,
 \label{eq:alpha-precoordenada-comun}
\end{equation}
avec \(p=\pi_{\mathrm{HMT}}-3\),
\(e_0=e_{\mathrm{HMT}}-2\) et
\(f=\varphi_{\mathrm{HMT}}-1\). De manière équivalente,
\(a=\pi_{\mathrm{HMT}}+e_{\mathrm{HMT}}-
\varphi_{\mathrm{HMT}}-4-\kappa_{\mathrm{per}}\). On n'introduit ici aucune valeur
conventionnelle de \(\alpha\) : \(p,e_0,f\) sont les trois caractères régionaux déjà
engendrés et \(\kappa_{\mathrm{per}}\) est la lecture rationnelle du registre
terminal. La clôture entière applique à
\eqref{eq:alpha-precoordenada-comun} la normalisation de
\eqref{eq:alpha-recurrencia} ; la coordinatisation qui suit exprime sur la même
précoordonnée une représentation analytique compatible. Son graphe de
dépendance reçoit \((p,e_0,f,\kappa_{\mathrm{per}})\) directement : il ne
contient pas d'arête \(\alpha_A\to\lambda\), bien que l'égalité démontrée plus
bas permette ensuite d'identifier les deux désignations de représentation.

Écrivons \(q=1/729\), soit \(F_{\mathsf s}\) le corps ordonné engendré par
les coordonnées de \(\mathsf s\), et posons
\(E_{9,\mathsf s}(X)=E_{21/22}^{(9)}(\mathsf s;X)\). Nous définissons le coefficient
de liaison
\begin{equation}
 \lambda(\mathsf s):=
 -\frac{E_{9,\mathsf s}(a(\mathsf s))
         \bigl(1-q\,a(\mathsf s)^3\bigr)}{a(\mathsf s)^{10}}.
 \label{eq:alpha-lambda-estado}
\end{equation}
Tous ses arguments appartiennent à l'état avant la reconnaissance
métrologique. En particulier, \(\lambda\) n'est pas obtenu en ajustant une chaîne
décimale cible. L'équation~\eqref{eq:alpha-lambda-estado} doit être lue selon
son type exact : une fois que la première représentation de l'état a produit
\(a(\mathsf s)\), elle fixe l'unique prolongement de la classe géométrique
définie ci-après. Elle ne constitue pas une seconde sélection indépendante de
\(a(\mathsf s)\), mais une représentation analytique corrélée du même
état enrichi.

En effet, pour \(\mu\in F_{\mathsf s}\), posons
\[
 E_{\mathsf s,\mu}(X)
 :=E_{9,\mathsf s}(X)+\mu\frac{X^{10}}{1-qX^3}.
\]
Comme \(0<a(\mathsf s)<10^{-2}\) et \(qa(\mathsf s)^3<1\), on a
\begin{equation}
 E_{\mathsf s,\mu}(a(\mathsf s))=0
 \quad\Longleftrightarrow\quad
 \mu=\lambda(\mathsf s).
 \label{eq:alpha-unicidad-lambda-clase-geometrica}
\end{equation}
Ainsi, le germe de coefficients ne contient pas de liberté une fois fixés l'état
et la classe de prolongements \(q\)-géométriques. Cette unicité est relative à la
classe déclarée ; elle n'affirme pas que le jet d'ordre neuf sélectionne par lui-même une
queue parmi toutes les séries analytiques possibles.

L'action locale sur des blocs de trois coefficients est
\begin{equation}
 \mathcal C_\alpha:F_{\mathsf s}^{3}\longrightarrow F_{\mathsf s}^{3},
 \qquad \mathcal C_\alpha(u,v,w)=q(u,v,w),
 \qquad c^{(0)}=(\lambda(\mathsf s),0,0).
 \label{eq:alpha-accion-coeficientes}
\end{equation}
Par conséquent, pour tout \(n\geq0\),
\begin{equation}
 b_{10+3n}=q^n\lambda(\mathsf s),\qquad
 b_{11+3n}=b_{12+3n}=0.
 \label{eq:alpha-recurrencia-coeficientes}
\end{equation}
C'est la règle qui calcule chaque coefficient au-delà du degré neuf. Il ne
reste aucun paramètre libre à chaque prolongation.

Pour \(M\geq0\), soit
\begin{equation}
 E_{\mathsf s,M}(X):=E_{9,\mathsf s}(X)
 +\lambda(\mathsf s)\sum_{n=0}^{M}q^nX^{10+3n}.
 \label{eq:alpha-truncamientos-completos}
\end{equation}
La fonction complète est
\begin{equation}
 E_{\mathsf s,\infty}(X)
 :=E_{9,\mathsf s}(X)
 +\lambda(\mathsf s)\frac{X^{10}}{1-X^3/729}.
 \label{eq:alpha-funcion-completa-explicita}
\end{equation}

Pour formuler le système projectif sans types implicites, nous fixons \(d_M=12+3M\) et définissons la fibration de jets
\begin{equation}
 \mathfrak J_M:=
 \coprod_{\mathsf s\in\mathcal X_\alpha^{\mathrm{enr}}}
 \{\mathsf s\}\times F_{\mathsf s}[X]/(X^{d_M+1}).
 \label{eq:alpha-espacio-jets-rev08}
\end{equation}
Si \(M'\ge M\), la restriction est
\begin{equation}
 \tau_{M',M}(\mathsf s,[P]_{d_{M'}})
 :=(\mathsf s,[P]_{d_M}).
 \label{eq:alpha-truncamiento-jets-rev08}
\end{equation}

\begin{proposition}[Naturalité de la récurrence analytique]
\label{prop:alpha-naturalidad-recurrencia-explicita}
Soit
\[
 \Gamma_M(\mathsf s)
 =\bigl(\mathsf s,[E_{\mathsf s,M}]_{d_M}\bigr).
\]
Si \(M'\geq M\), la troncature \(\tau_{M',M}\) satisfait
\begin{equation}
 \tau_{M',M}\circ\Gamma_{M'}=\Gamma_M,
 \qquad
 \Gamma_{E21}(\mathsf s)=\varprojlim_M\Gamma_M(\mathsf s).
 \label{eq:alpha-naturalidad-explicita}
\end{equation}
La projection sur les trois nouveaux degrés est précisément
\(\mathcal C_\alpha^{M+1}(c^{(0)})\).
\end{proposition}

\begin{proof}
L'équation~\eqref{eq:alpha-recurrencia-coeficientes} fixe le bloc de degrés \(10+3n,11+3n,12+3n\). Lors du passage de \(M\) à \(M+1\), seul le bloc suivant est ajouté ; le tronquer redonne littéralement le polynôme précédent. La compatibilité pour \(M'\geq M\) s'obtient par composition. Ainsi, la limite inverse n'est pas une section choisie du jet nu, mais l'orbite unique de \(c^{(0)}\) sous \(\mathcal C_\alpha\).
\end{proof}

\subsubsection{Convergence uniforme, racine simple et intervalles rationnels}
\label{subsec:alpha-cotas-completas-rev08}

Les cylindres régionaux et le registre périodique fournissent des bornes au moyen
de l'arithmétique des intervalles rationnels. Pour chaque coordonnée
\[
 c\in\{\pi_{\mathrm{HMT}},e_{\mathrm{HMT}},\varphi_{\mathrm{HMT}}\},
\]
soit \(P_c\) son préfixe nonadique certifié de mille chiffres. La donnée utilisée
n'est pas le rationnel \(P_c\) comme s'il s'agissait de la valeur complète, mais le cylindre
semi-ouvert et sa fermeture extérieure
\[
 \mathcal C_c^{(1000)}=[P_c,P_c+\varepsilon),
 \qquad
 \overline{\mathcal C}_c^{(1000)}=[P_c,P_c+\varepsilon],
 \qquad \varepsilon=10^{-1000}.
\]
L'arithmétique dirigée opère de manière conservative avec ces fermetures. Si
\(P_\pi,P_e,P_\varphi\) désignent les trois préfixes, on obtient
\begin{equation}
\begin{aligned}
 a_-&:=P_\pi+P_e-(P_\varphi+\varepsilon)-4-
       \kappa_{\mathrm{per}},\\
 a_+&:=(P_\pi+\varepsilon)+(P_e+\varepsilon)-P_\varphi-4-
       \kappa_{\mathrm{per}},\\
 A_{1000}&:=[a_-,a_+],
 \qquad a(\mathsf s)\in[a_-,a_+]=A_{1000},
 \qquad \operatorname{diam}A_{1000}=3\varepsilon.
\end{aligned}
\label{eq:alpha-propagacion-colas-rev08}
\end{equation}
L'extension naturelle par intervalles de
\eqref{eq:alpha-lambda-estado} transporte ensuite \(A_{1000}\), le cylindre
de \(\pi_{\mathrm{HMT}}\) et la borne rationnelle dirigée de \(\delta\) vers un
intervalle fermé \(\Lambda_{1000}\ni\lambda(\mathsf s)\) ; aucune
des trois queues n'est figée. En particulier,
\begin{equation}
\begin{aligned}
 \frac{7297352569283800997285105472380662}{10^{36}}
 &<a(\mathsf s)<
 \frac{7297352569283800997285105472380663}{10^{36}},\\
 -8.711\mathbin{\times}10^{-6}
 &<\lambda(\mathsf s)<-8.709\mathbin{\times}10^{-6}.
\end{aligned}
 \label{eq:alpha-cotas-precoordenada-lambda}
\end{equation}
Ces inégalités utilisent des préfixes certifiés avec leur borne rationnelle de
queue ; elles ne reçoivent aucune table d'alpha. Le vérificateur du paquet les reproduit
à partir des trois sorties nonadiques et de la coordinatisation terminale.

Soit
\[
 r:=\frac{10^{-6}}{729}<1.372\mathbin{\times}10^{-9}.
\]
Pour \(0\leq x\leq10^{-2}\), la queue postérieure à \(M\) satisfait
\begin{equation}
 \left|E_{\mathsf s,\infty}(x)-E_{\mathsf s,M}(x)\right|
 \leq
 8.711\mathbin{\times}10^{-26}\,\frac{r^{M+1}}{1-r}.
 \label{eq:alpha-cota-cola-uniforme}
\end{equation}
De plus,
\begin{equation}
 \left|\frac{\mathrm d}{\mathrm dx}
  \left(\lambda\frac{x^{10}}{1-x^3/729}\right)\right|
 \leq 8.711\mathbin{\times}10^{-6}
 \left(\frac{10\,10^{-18}}{1-r}
 +\frac{3\,10^{-24}/729}{(1-r)^2}\right)
 <8.712\mathbin{\times}10^{-23}.
 \label{eq:alpha-cota-derivada-cola}
\end{equation}
La queue des dérivées postérieure au bloc \(M\) admet en outre l'estimation dépendant de la profondeur
\begin{equation}
 \sup_{0\le x\le10^{-2}}
 \left|E'_{\mathsf s,\infty}(x)-E'_{\mathsf s,M}(x)\right|
 \le 8.711\mathbin{\times}10^{-24}\,r^{M+1}
 \left(\frac{13+3M}{1-r}+\frac{3r}{(1-r)^2}\right)
 \xrightarrow[M\to\infty]{}0.
 \label{eq:alpha-cota-resto-derivadas-rev08}
\end{equation}

\begin{theorem}[Fonction limite et unique racine simple de la coordinatisation corrélée]
\label{thm:alpha-raiz-completa-explicita}
La suite \(E_{\mathsf s,M}\) converge uniformément, ainsi que ses
dérivées, vers \(E_{\mathsf s,\infty}\) sur
\(\overline I_\alpha=[0,10^{-2}]\). On a
\begin{equation}
 E_{\mathsf s,\infty}(a(\mathsf s))=0,
 \qquad
 \inf_{x\in\overline I_\alpha}
 E'_{\mathsf s,\infty}(x)>6.239.
 \label{eq:alpha-raiz-y-derivada-completas}
\end{equation}
Par conséquent, \(a(\mathsf s)\) est l'unique racine simple de la fonction complète dans \(I_\alpha\).
\end{theorem}

\begin{proof}
Le rapport de deux termes consécutifs de la queue est borné par \(r<1\),
ce qui donne \eqref{eq:alpha-cota-cola-uniforme} ; la dérivation terme à
terme et la somme des deux séries géométriques donnent
\eqref{eq:alpha-cota-derivada-cola} et
\eqref{eq:alpha-cota-resto-derivadas-rev08}. La preuve du jet établit
\(E'_{9,\mathsf s}>6.24\) dans \(\overline I_\alpha\) ; par conséquent, la
perturbation complète maintient la dérivée au-dessus de \(6.239\).
Enfin, en substituant \eqref{eq:alpha-lambda-estado} dans
\eqref{eq:alpha-funcion-completa-explicita}, les deux termes évalués en
\(a(\mathsf s)\) s'annulent exactement. La fonction est strictement
croissante ; son unique racine est donc simple et coïncide avec cette précoordonnée.
\end{proof}

Pour rendre effective l'identification à toute profondeur, les évaluations
dirigées des extrémités certifient d'abord
\[
 E_{\mathsf s,\infty}(0)<0<
 E_{\mathsf s,\infty}(10^{-2}),
 \qquad
 E'_{\mathsf s,\infty}\geq\mu:=6239/1000>c:=6
 \quad\hbox{en }[0,10^{-2}].
\]
Par continuité, il existe une racine dans l'intervalle et, par la deuxième inégalité, elle est unique. Ce n'est qu'après avoir établi cette existence, la monotonie et l'invariant selon lequel l'encadrement contient la racine que l'on applique le rattrapage par le théorème des accroissements finis. Nous partons de \(D_0=[0,10^{-2}]\). Si \(D_j=[r_j,s_j]\), soit \(m_j=(r_j+s_j)/2\), et soit \([F_j^-,F_j^+]\) une inclusion rationnelle dirigée de \(E_{\mathsf s,\infty}(m_j)\), obtenue en propageant les cylindres de \(\pi_{\mathrm{HMT}},e_{\mathrm{HMT}},\varphi_{\mathrm{HMT}}\), \(\delta\) et \(\lambda\), avec
\begin{equation}
 0\leq F_j^+-F_j^-\leq \frac{c}{4}\,(s_j-r_j).
 \label{eq:alpha-anchura-evaluacion-rescate-rev08}
\end{equation}
Lorsque \(F_j^+<0\), on conserve la moitié droite ;
lorsque \(F_j^->0\), la moitié gauche. Si
\(F_j^-\leq0\leq F_j^+\), on ne cherche pas à décider si le milieu est
exactement la racine. Soit \(w_j=s_j-r_j\). La borne uniforme
\(E'_{\mathsf s,\infty}\geq c=6\) et le théorème des accroissements finis impliquent que
toute racine contenue dans \(D_j\) appartient à
\begin{equation}
 D_{j+1}=D_j\cap
 \left[m_j-\frac{w_j}{4},
             m_j+\frac{w_j}{4}\right].
 \label{eq:alpha-rescate-derivada-rev08}
\end{equation}
En effet, soit \(\xi_j\) la racine et soit
\(y_j=E_{\mathsf s,\infty}(m_j)\in[F_j^-,F_j^+]\). Comme l'encadrement contient
zéro et satisfait \eqref{eq:alpha-anchura-evaluacion-rescate-rev08}, on a
\(|y_j|\leq cw_j/4\). Le théorème des accroissements finis fournit
\(|m_j-\xi_j|\leq |y_j|/c\leq w_j/4\), ce qui prouve
\eqref{eq:alpha-rescate-derivada-rev08}. L'argument inclut le cas
\(y_j=0\) : la racine est alors au milieu, mais l'algorithme n'a pas
besoin de reconnaître cette égalité pour la conserver.

Si l'encadrement disponible ne satisfait pas \eqref{eq:alpha-anchura-evaluacion-rescate-rev08}, on raffine d'abord les cylindres sources et on répète l'évaluation ; on ne déclare aucun signe et on n'accepte pas de contraction insuffisante. Dans chacune des trois branches, le nouvel intervalle conserve la racine et a un diamètre au plus égal à \(w_j/2\). Pour chaque profondeur cible finie, le raffinement dirigé permet de satisfaire la borne de largeur sans comparer de manière non certifiée un nombre réel à zéro. La monotonie démontre par récurrence
\begin{equation}
 a(\mathsf s)\in D_{j+1}\subseteq D_j,\qquad
 \operatorname{diam}D_{j+1}\leq
 \tfrac12\operatorname{diam}D_j,
 \qquad \operatorname{diam}D_j\longrightarrow0,
 \label{eq:alpha-biseccion-racional}
\end{equation}
où l'inégalité est une borne de contraction et non une égalité imposée.
Le certificat exécutable applique
cette règle à vingt-quatre niveaux en base mille et vérifie que la clôture
extérieure complète \(A_{1000}\), et non sa seule extrémité inférieure, est contenue dans
chaque cylindre exprimé. Ses tests négatifs rejettent un encadrement trop
large, un encadrement désordonné, un point autre que le milieu, un domaine sans
changement de signe et une fourchette dégénérée. Le prolongement coinductif permet
de remplacer mille par toute profondeur requise.

Concrètement, pour \(S_N=1000^N\), le vérificateur calcule
\begin{equation}
 j_N^-:=\lfloor S_Na_-\rfloor,
 \qquad j_N^+:=\lfloor S_Na_+\rfloor
 \label{eq:alpha-indices-cilindro-dirigidos-rev08}
\end{equation}
et n'exprime le niveau que si \(j_N^-=j_N^+\). Cette égalité, qui inclut
l'extrémité supérieure de la clôture extérieure, prouve
\(A_{1000}\subset[j_N^-/S_N,(j_N^-+1)/S_N)\) et évite d'attribuer à la cellule
gauche une extrémité située exactement sur sa frontière droite.

Soit \(C_N^A(\mathsf s)\) le cylindre entier de \eqref{eq:alpha-cilindros-via-a-exactos}. Nous choisissons récursivement \(m(N)>m(N-1)\) de sorte que \(\operatorname{diam}D_{m(N)}<1000^{-N}/4\), et définissons
\begin{equation}
I_{B,N}:=D_{m(N)}\cap C_N^A(\mathsf s).
 \label{eq:alpha-intervalos-b-completos}
\end{equation}
Nous définissons
\(\ell_N:=\inf I_{B,N}\) et \(u_N:=\sup I_{B,N}\). Tous deux sont rationnels ;
l'intervalle est non vide parce que les deux facteurs contiennent
\(a(\mathsf s)\). Le cylindre étant semi-ouvert, \(u_N\) peut ne pas
appartenir à \(I_{B,N}\). L'évaluation en cette extrémité s'entend au moyen
du prolongement continu de \(E_{\mathsf s,\infty}\) à la fermeture extérieure. Ainsi,
\begin{equation}
 I_{B,N+1}\subseteq I_{B,N}\subseteq C_N^A(\mathsf s),
 \qquad \operatorname{diam}I_{B,N}\longrightarrow0,
 \qquad E_{\mathsf s,\infty}(\ell_N)\leq0\leq
 E_{\mathsf s,\infty}(u_N).
 \label{eq:alpha-inclusion-signos-completa}
\end{equation}
La borne~\eqref{eq:alpha-cota-cola-uniforme} permet de choisir une
troncature qui conserve le signe à chaque extrémité où la fonction complète
ne s'annule pas. Si une extrémité coïncide avec la racine, on utilise l'identité de
la fonction complète et l'appartenance déjà démontrée ; on n'attribue pas ce même
signe à ses troncatures. En effet, le reste exact est
\[
 E_{\mathsf s,\infty}(x)-E_{\mathsf s,M}(x)
 =\lambda(\mathsf s)x^{10}
   \frac{(qx^3)^{M+1}}{1-qx^3}.
\]
Dans le domaine positif considéré, où \(\lambda(\mathsf s)<0\), on obtient à la racine
\[
 E_{\mathsf s,M}(a(\mathsf s))
 =-\lambda(\mathsf s)a(\mathsf s)^{10}
   \frac{(qa(\mathsf s)^3)^{M+1}}{1-qa(\mathsf s)^3}>0
 \qquad(M<\infty).
\]
La certification des signes au sens large dans \eqref{eq:alpha-inclusion-signos-completa} se rapporte donc à \(E_{\mathsf s,\infty}\). Son évaluation utilise la fonction complète ou la troncature accompagnée du reste signé. Cela couvre aussi bien l'infimum appartenant au cylindre que le supremum éventuellement exclu de celui-ci.

\begin{proposition}[Carrés de compatibilité et cylindres communs]
\label{prop:alpha-cuadrados-compatibilidad}
Les restrictions de mots et de représentations analytiques commutent :
\[
\begin{array}{ccc}
 \mathcal X_\alpha^{\mathrm{enr}}&
 \xrightarrow{\ \mathsf T_{\alpha,N'}\ }&\mathcal W_{N'}\\[1mm]
 \Big\Vert&&\Big\downarrow\rho_{N',N}\\[1mm]
 \mathcal X_\alpha^{\mathrm{enr}}&
 \xrightarrow{\ \mathsf T_{\alpha,N}\ }&\mathcal W_N
\end{array}
\qquad
\begin{array}{ccc}
 \mathcal X_\alpha^{\mathrm{enr}}&
 \xrightarrow{\ \Gamma_{M'}\ }&\mathfrak J_{M'}\\[1mm]
 \Big\Vert&&\Big\downarrow\tau_{M',M}\\[1mm]
 \mathcal X_\alpha^{\mathrm{enr}}&
 \xrightarrow{\ \Gamma_M\ }&\mathfrak J_M .
\end{array}
\]
La collection \(I_{B,N}\) assure la connexion effective entre les deux
carrés et satisfait matériellement
\(I_{B,N}\subseteq C_N^A\) pour tout \(N\).
\end{proposition}

\begin{proof}
Le premier carré est la compatibilité de la normalisation entière ; le second est la proposition~\ref{prop:alpha-naturalidad-recurrencia-explicita}. L'inclusion découle de la définition~\eqref{eq:alpha-intervalos-b-completos}, et la non-vacuité et les signes découlent du théorème~\ref{thm:alpha-raiz-completa-explicita}. Aucune étape n'identifie la racine d'une troncature à celle d'une queue arbitraire.
\end{proof}

Nous définissons maintenant, sans laisser implicite le lecteur de préfixes,
\begin{equation}
 \alpha_B(\mathsf s):=
 \operatorname*{arg\,zero}_{x\in I_\alpha}E_{\mathsf s,\infty}(x),
 \qquad
 \operatorname{pref}_N(x):=
 \frac{\lfloor1000^Nx\rfloor}{1000^N},
 \qquad
 \alpha_{B,N}:=\operatorname{pref}_N(\alpha_B).
 \label{eq:alpha-definicion-b-y-prefijos-rev08}
\end{equation}

\begin{theorem}[Égalité projective des deux représentations corrélées]
\label{prop:alpha-dos-limites}
Pour tout \(N\), la voie entière et la voie analytique expriment le même préfixe :
\begin{equation}
 \boxed{\alpha_{A,N}=\operatorname{pref}_N(\alpha_A)
 =\operatorname{pref}_N(\alpha_B)=\alpha_{B,N}.}
 \label{eq:alpha-dos-vias-cada-nivel}
\end{equation}
À la limite,
\begin{equation}
 \boxed{\alpha_A
 =\operatorname*{arg\,zero}_{x\in I_\alpha}
       E_{\mathsf s,\infty}(x)
 =\alpha_B=\alpha_{\mathrm{HMT}}.}
 \label{eq:alpha-dos-vias-limite-coinductivo}
\end{equation}
\end{theorem}

\begin{proof}
La normalisation entière exprime exactement \(a(\mathsf s)\). Le
théorème~\ref{thm:alpha-raiz-completa-explicita} démontre que la coordinatisation
analytique représente ce même point comme unique racine simple. Les
cylindres semi-ouverts \(C_N^A\) fixent un préfixe canonique et
\(I_{B,N}\subseteq C_N^A\) contient \(\alpha_B=\alpha_A\). Leurs
préfixes coïncident donc pour chaque \(N\) ; l'emboîtement et les diamètres tendant vers
zéro donnent l'égalité des limites. La reconnaissance
métrologique intervient ensuite.
\end{proof}

La fonction analytique constitue une représentation corrélée du même
état, non une seconde dérivation causalement indépendante de la clôture
entière. L’égalité identifie le même point au moyen de la clôture entière
et de la représentation analytique corrélée, toutes deux construites à partir
de l’état enrichi. Le prolongement est déterminé par la
précoordonnée et la récurrence explicite, non par le seul jet d’ordre neuf ;
la reconnaissance conventionnelle est postérieure à la génération.

